\section{Resumen}
En este informe se realizaron mediciones y ensayos experimentales sobre un circuito con acoplamiento magnético utilizando un transformador de inductancias primaria ($L_1$) y secundaria ($L_2$) inicialmente desconocidas, operando a una frecuencia de trabajo de $50\text{ kHz}$. A partir de las caracterizaciones efectuadas mediante el analizador de impedancias, se determinaron los parámetros reales de los bobinados y se estimó el coeficiente de acoplamiento magnético ($K$). 

Durante la experiencia, se analizó el comportamiento del dispositivo al ser configurado tanto en modo elevador como reductor de tensión, y se comprobó el efecto de aislamiento eléctrico del transformador ante variaciones de un offset de continua en la señal de entrada. Asimismo, se estudió cuantitativamente el efecto de cargar el circuito mediante la incorporación de una resistencia de carga ($R_L$), evaluando la impedancia de entrada ($Z_e$) en ambos sentidos de conexión del transformador con el propósito de optimizar la transferencia energética y evitar sobrecargar el generador de señales.

En general, los resultados experimentales obtenidos presentaron una buena concordancia con los modelos teóricos de transformadores reales y las simulaciones asociadas, permitiendo verificar de manera práctica la influencia de las inductancias de dispersión, las resistencias de los devanados y las pérdidas en el núcleo sobre las características eléctricas globales del sistema.

\newpage
